\documentclass[10pt,a4paper]{amsart}
\usepackage[margin=19mm]{geometry}
\usepackage{amsmath,amssymb,mathtools}
\usepackage[scr=rsfs,cal=euler]{mathalfa}
\usepackage{xfrac}
\usepackage[unicode,hidelinks,bookmarks=false]{hyperref}
\newcommand{\ie}{i.e.\ }
\newcommand{\cf}{cf.\ }
\newcommand{\resp}{respectively\ }
\newcommand{\Subsection}{\subsection}
\providecommand{\nobreakdash}{\nobreak\hbox{-}}
\setlength{\emergencystretch}{2em}
\multlinegap0pt
\usepackage{algorithm}\usepackage{algpseudocode}
\usepackage{mathtools}
\usepackage{float}
\usepackage{xcolor}
\usepackage[shortlabels]{enumitem}
\usepackage{tcolorbox}
\usepackage{amssymb}
\newtheorem{lemma}{Lemma}
\newcommand{\RR}{\mathbb{R}}
\newcommand{\Var}{\mathbb{V}}
\newcommand{\inlaw}{\buildrel {\mathcal L} \over =}
\newcommand{\isdef}{\buildrel {\rm def} \over =}
\title{The Pearson IV distribution: Random variate generation and applications}
\author{Luc Devroye, Joe R. Hill}
\date{Source: arXiv:2605.01586v1 (CC BY 4.0)}
\begin{document}
\maketitle

\noindent\emph{Source and licence.} The Pearson IV distribution: Random variate generation and applications. Version arXiv:2605.01586v1 (CC BY 4.0), distributed by arXiv under the Creative Commons Attribution 4.0 International licence, http://creativecommons.org/licenses/by/4.0/, as recorded in the arXiv licence metadata for this exact version and shown on the abstract page https://arxiv.org/abs/2605.01586v1. Full licence notice: \texttt{licenses/arxiv-2605.01586-CC-BY-4.0.txt}.\par\smallskip

\noindent\textit{Editorial scope.} Counted unit: the article's lemma lemma (source0.tex lines 741--747). The article's complete original proof of it is delivered with it (source0.tex lines 749--782, one \texttt{proof} environment of 34 lines). No mathematics is written, repaired or re-derived: the statement and the proof are the article's own lines, in the article's own notation, and no retained line is substituted or re-flowed.  Retained uncounted as the source-local context the proof needs: lines 420--426; lines 428--433; lines 442--446; lines 462--464; lines 697--733; lines 790--812.  Nothing was omitted from any retained span, so the package carries no omission record.  No retained line contains a citation, so the wrapper carries no bibliography and no bibliography entry.  No retained line contains a figure environment, an image-inclusion command or drawing code, so no image is delivered and the package declares no asset dependency.  Editorial formatting only, in addition: an AMS \texttt{amsart} preamble that restates the article's own declarations and macros used by the retained lines, namely the environments \texttt{lemma} on separate counters, together with the standard packages those lines need.  The retained environments print in this document as Lemma 1 in order of appearance, whereas the article prints the same items with the numbering of its own full document.  The complete native source of the article as distributed in the arXiv e-print for this exact version is bundled unchanged as \texttt{sources/199647/source0.tex}.
\begin{equation}\label{atan}
h(y) =  
\begin{cases}
\gamma e^{sy} ( \cos^2 (y) )^{a-1} & \text{if} ~|y| \le \frac{\pi }{ 2}, \\
0 & \text{else.} 
\end{cases}
\end{equation}

\begin{align}\label{derivatives}
g' (y) &= s - 2(a-1) \tan (y) ,\\
g'' (y) &= - \frac{2 (a-1)} {\cos^2 (y)}, \\
g''' (y) &= - \frac{4 (a-1) \tan (y)} {\cos^2 (y)} ,\\
g'''' (y) &= - \frac{4 (a-1) ( 1 + 2 \sin^2 (y)) } {\cos^4 (y)}.
\end{align}

\begin{align}\label{variance}
\Var \{ X \} 
= \frac{1}{4} \left( \psi_1 (a+ is/2) + \psi_1 (a- is/2) \right)
\in \left[ \frac{2a}{s^2 + 4a^2} , \frac{2(a+1)}{s^2 + 4a^2}\right]
\end{align}

\begin{align}\label{fradelizi}
\frac{1}{\sqrt{12} } \le h(m)\sqrt{\Var\{X\}} \le \frac{1}{\sqrt{2}}\,.
\end{align}

We can derive upper bounds for the arctan-mapped density $h$ (see \eqref{atan}) based on
the derivatives of $g=\log h$ (see \eqref{derivatives}) and standard Taylor series bounds:
\begin{align}
g(y) 
&\le 
\begin{cases}
g(m) + \frac{(y-m)^2}{2} g''(m), & y \ge m , \cr
g(m) + \frac{(y-m)^2}{2} g''(m) - \frac{(y-m)^3}{6} g'''(m), & y \le m, \cr
\end{cases}  \\
&\le 
\begin{cases}
g(m) + \frac{(y-m)^2}{2} g''(m), & y \ge m , \cr
g(m) + \frac{(y-m)^2}{4} g''(m),  & m - D \le y \le m, \cr
\end{cases}\label{gaussian2}
\end{align}
whenever 
$$
\frac {D |g'''(m)|}{6} \le \frac{ |g''(m)|}{4}.
$$
This is satisfied if we pick $D$ such that $D \beta \le 3/4$.
For example, if we set $D=\pi$, then for all $m$ satisfying
\begin{align}\label{condition}
\tan (m) \le \frac{3}{4 \pi} \approx 0.2387324,
\end{align}
we have
\begin{align}\label{gaussian3}
h(y) \le 
% h(m) e^{-(a-1)(1+\beta^2) \frac{(y-m)^2}{2}}
h(m) e^{-\tau^2 \frac{(y-m)^2}{2}}, y \in \RR,
\end{align}
where
\begin{align*}
\tau \isdef \sqrt {|g''(m)|/2} = \sqrt{(a-1)(1+\beta^2)}.
\end{align*}
A condition equivalent to \eqref{condition}
is $|m| \le \arctan \left( \frac{3}{4\pi} \right) \approx 0.2344$. 
Algorithm  \ref{simplegaussian} uses rejection either from the Gaussian implied by \eqref{gaussian3} or rejection from the uniform density on $[-\pi/2,\pi/2]$.

\begin{algorithm}[H]
\caption{PearsonIV generator for $a > 1$ by rejection from the normal density when $\tan (m) \le 3/(4 \pi) \approx 0.2387324$, or equivalently, when $m \le  \arctan(3/(4 \pi)) \approx 0.2344$.}\label{simplegaussian}
\begin{algorithmic}[1]
\State let $\beta = s/(2(a-1))$,  $m = \arctan(\beta)$, $\tau = \sqrt{(a-1)(1+\beta^2 )}$
\If{$\tau \ge \sqrt{2/\pi}$}
\Repeat 
\State let $N$ be a standard Gaussian random variable 
\State let $U$ be uniform on $[0,1]$
\State set $Y \leftarrow  m + N/\tau$
\Until {$U h(m) e^{-\tau^2  \frac{(Y-m)^2}{2}}  \le h(Y)$} 
\State
\Comment{the acceptance condition is equivalent to $U e^{-N^2 /2} \le h(Y)/h(m)$}
\Else
\Repeat 
\State let $Y$ be uniform on $[-\pi/2, \pi/2]$ 
\State let $U$ be uniform on $[0,1]$
\Until {$U   \le h(Y)/h(m)$} 
\EndIf
\State \textbf {return} $X \leftarrow \tan(Y)$  
\Comment{$X \inlaw P_{a,s}$ }
\end{algorithmic}
\end{algorithm}
%------------------------------------------------------------------------------

\begin{lemma} 
The expected number of iterations taken by algorithm \ref{simplegaussian} is not more than
$$
\sqrt{2 \pi + \pi^2}  \approx 4.02,
$$
uniformly over all values of the parameters with $a \ge 1$ and $|m| \le \arctan \left( \frac{3}{4\pi} \right) \approx 0.2344$.
\end{lemma}

\begin{proof}
If we were to use \eqref{gaussian2}, then the  expected number of iterations before halting would be
$$ 
\frac{ \sqrt{2 \pi}  h(m) } { \tau  } .
$$
However, if we bound $h$ by $h(m)$ and use rejection from the uniform density on $[-\pi/2,\pi/2]$, then the expected number of iterations before halting
would be $\pi h(m)$.  Taking  the best of both leads to the cut-off
value at $\tau = \sqrt{2/\pi }$. 
Using \eqref{variance} and Fradelizi's inequality \eqref{fradelizi}, we have
\begin{align*}
h(m) \le \sqrt{ a + \frac{s^2}{4a}} .
\end{align*}
Assume first $\tau \le \sqrt{2/\pi }$. Then 
$a-1 \le 2/\pi$ and $s^2/(4(a-1)) \le 2/\pi$, which implies that
$$
\pi h(m) \le \pi \sqrt{1 + 4/\pi} = \sqrt{\pi^2 + 4}.
$$
If, on the other hand, $\tau \ge \sqrt{2/\pi }$, then
\begin{align*}
\frac{ \sqrt{2 \pi}  h(m) } { \tau  }
&\le
\frac{ \sqrt{2 \pi}  \sqrt{a +\frac{s^2}{4a} } } { \tau  }
\le
\frac{ \sqrt{2 \pi}  \sqrt{1 +(a-1) +\frac{s^2}{4(a-1)} } } { \tau  }\\
&=
\frac{ \sqrt{2 \pi}  \sqrt{1 + \tau^2 } } { \tau  }
=
\sqrt{2 \pi}  \sqrt{1 + \frac{1}{\tau^2} } 
\le 
\sqrt{2 \pi}  \sqrt{1 + \frac{\pi}{2} } 
=
\sqrt{2 \pi + \pi^2} .
\end{align*}
\end{proof}

\end{document}
